
\subsubsection{3.1 Framework Overview}

The BCBHS framework was designed around a two-stage verification structure. The first stage establishes that domain-specific health signals exist with sufficient discriminative power (H-E1) and that submission accumulation Granger-causes score compression in real panel data (H-M1). The second stage, contingent on the first, was intended to build a Cox proportional hazards survival model predicting time-to-discriminative-collapse and to validate prospective early-warning lead times. The second stage (H-M3, H-M4) was not executed in this work.

\subsubsection{3.2 Panel Construction}

\textbf{Data source.} The quarterly leaderboard panel was constructed from the \texttt{pwc-archive/evaluation-tables} HuggingFace archive, which provides historical Papers With Code leaderboard data from 2018 to 2025. The raw dataset contains 48,311 submission rows across 1,120 tasks.

\textbf{Qualifying benchmarks.} Benchmarks were filtered to those with at least 20 model submissions and at least 2 years of submission history. This yielded 466 qualifying benchmarks across CV and NLP domains.

\textbf{Panel structure.} Submissions were aggregated to quarterly resolution. For each benchmark-quarter pair (B, t), two derived quantities were computed: \texttt{score_var_top10} (variance of scores among the top-10 submitted models in that quarter) and \texttt{cumulative_count} (total submissions to that benchmark as of quarter t). The panel contains 6,938 observations.

\subsubsection{3.3 Domain-Specific Health Estimators H_d(B, t)}

The BCBHS framework uses domain-specific signals because CV, NLP, and tabular benchmarks exhibit different saturation phenotypes:

\textbf{CV — Score variance proxy.}

$$H_d^{\text{CV}}(B, t) = \text{score\_var\_top10}(B, t)$$

CV saturation was hypothesized to manifest as score \textit{convergence}: as models overfit the test distribution, their scores cluster, reducing variance among the top-k. A lower H_d^CV value signals saturation. Empirically, this direction was confirmed (see Section 5.1).

\textbf{NLP — Contamination-adjusted deviation signal.}

$$H_d^{\text{NLP}}(B, t) = \text{NMD}(B, t)$$

NLP saturation was hypothesized to manifest as score \textit{divergence}: contaminated models produce anomalously high scores, increasing the normalized mean deviation (NMD) of the top-k score distribution. A higher H_d^NLP signals saturation. The originally planned ConStat S_index was unavailable (PWC REST API shut down July 2025); NMD was used as a fallback.

\textbf{Tabular — Block-bootstrapped Kendall τ rank stability.}

$$H_d^{\text{Tab}}(B, t) = \hat{\tau}_{\text{block}}(B, t)$$

Tabular saturation was hypothesized to manifest as premature rank stabilization. The block-bootstrapped Kendall rank correlation between top-k orderings at quarters t and t−1 was computed with 1,000 bootstrap iterations over model family blocks. A higher $\hat{\tau}_{\text{block}}$ signals saturation.

All three signals use a 24-month (8-quarter) lookback window, selected based on discriminability analysis across multiple lookback window lengths (Section 5.1).

\subsubsection{3.4 Compression Event Detection}

A compression event for benchmark B at quarter t is defined as:

$$\text{compression\_event}(B, t) = \mathbf{1}\left[\text{score\_var\_top10}(B, t) < \mu_\sigma - 1.5 \cdot \hat{\sigma}_{\text{measurement}}(B)\right]$$

where $\hat{\sigma}_{\text{measurement}}(B)$ is estimated from repeated submissions of the same model within a quarter. The event requires the threshold to be sustained for at least 2 consecutive quarters. The median $\hat{\sigma}_{\text{measurement}}$ across 7,592 benchmarks (from the broader pre-filter archive) is 0.3323. The broader benchmark set was used for sigma estimation to avoid overfitting the noise threshold to the smaller 466-benchmark qualified subset.

\subsubsection{3.5 Granger Causality Analysis}

To test whether submission accumulation temporally precedes score compression, Granger causality analysis was applied to the quarterly panel. For each benchmark with at least 9 quarterly observations (lag + 5 minimum), a vector autoregression (VAR) model was estimated with \texttt{cumulative_count} and \texttt{score_var_top10} as endogenous variables, following ADF stationarity testing with iterative first-differencing. Both forward (submissions → compression) and reverse (compression → submissions) directions were tested. The panel-level result reports the minimum p-value across all benchmarks with sufficient time series length.

Granger causality establishes temporal predictability — that past values of submission count improve forecasts of score compression beyond the latter's own history — but does not rule out confounding variables such as benchmark age or task adoption lifecycle.

\subsubsection{3.6 Existence Validation Approach (H-E1)}

To validate H_d discriminative power in controlled conditions prior to real-data analysis, synthetic panels of 20 saturated and 20 healthy benchmarks per domain were constructed, calibrated to match real PWC panel statistics (median $\hat{\sigma}_{\text{measurement}} = 0.3323$, empirical submission count distribution). The gate criterion was Mann-Whitney U p < 0.05 AND |Cohen's d| > 0.5 in at least 2 of 3 domains.

\textbf{Pipeline summary:}

```
PWC HuggingFace archive (48,311 rows)
    → filter (≥20 submissions, ≥2 years)
Quarterly panel (6,938 rows × 466 benchmarks)
    → domain-specific H_d computation
H_d signals (CV: score_var_top10; NLP: NMD; Tabular: τ_block)
    → compression detection (1.5σ, ≥2 consecutive quarters)
Compression events (389 events, 145 benchmarks, 31.1\%)
    → Granger causality (ADF stationarity → VAR → F-test)
Granger result: submissions → compression (p=1.854e-05, lag=2)
    → [NOT EXECUTED: collapse event recalibration → Cox PH model]
```

